\section{Problem Summary}
Generate all binary strings of length \(n\) in an order where every consecutive pair differs in exactly one bit.

\section{Editorial}
A standard closed-form construction exists for Gray code:
\[
g(i) = i \oplus (i \gg 1),
\]
where \(i\) runs from \(0\) to \(2^n - 1\).

Why does this help? Consecutive integers differ at one lowest changed bit plus a block of lower bits, and the xor-with-shift transformation folds that pattern into a representation where only one output bit changes between neighbors.

So we can:
\begin{itemize}
\item iterate \(i\) from \(0\) to \(2^n - 1\)
\item compute \(g(i)\)
\item print \(g(i)\) as a binary string of exactly \(n\) bits
\end{itemize}

This is much simpler than building the list recursively, although the recursive reflection construction would also be valid.

\section{Pseudocode}
Set the limit to \(2^n\).

For each integer \(i\) from \(0\) to \(\text{limit} - 1\):
\begin{itemize}
\item compute \(g = i \oplus (i \gg 1)\)
\item print the \(n\)-bit binary representation of \(g\)
\end{itemize}

\section{Complexity Analysis}
\begin{itemize}
\item Time: \(O(n \cdot 2^n)\), dominated by printing \(2^n\) strings of length \(n\)
\item Memory: \(O(1)\) besides the output itself
\end{itemize}

\section{Notes}
Make sure to print leading zeroes so every line has exactly \(n\) characters.
